\chapter{Conclusão}
\label{conclusao}

Portanto, o construção desse relatório permitiu consolidar a proposta conceitual do projeto Micromouse, ao reunir as definições iniciais das áreas de Estrutura, Energia, Eletrônica e Software. Ao longo do documento, foram apresentados os objetivos do projeto, os requisitos principais, a organização da equipe, o dimensionamento mecânico, a análise energética, a arquitetura do sistema embarcado, o planejamento de software, o orçamento e o cronograma de execução.

Com base nas análises realizadas, conclui-se que a solução proposta é viável para a construção de um robô móvel autônomo capaz de atuar em labirintos de diferentes dimensões, respeitando as restrições físicas, energéticas e funcionais estabelecidas pela disciplina. A escolha dos componentes, como a ESP32, os sensores ToF, os motores N20 com encoder e a bateria Li-Ion, mostrou-se compatível com as necessidades do sistema.

Assim, esta etapa do projeto fornece uma base técnica e organizacional para as próximas fases, que envolvem a construção do protótipo, a integração entre os subsistemas, a implementação do firmware e do dashboard, além da realização de testes práticos para validação do funcionamento do Micromouse.
